% ============================================================
%  Projektdokumentation – Informatik 3
% ============================================================
\documentclass[12pt, a4paper, twoside]{article}

% ── Encoding & Sprache ──────────────────────────────────────
\usepackage[utf8]{inputenc}
\usepackage[T1]{fontenc}
\usepackage[ngerman]{babel}

% ── Schrift ─────────────────────────────────────────────────
\usepackage{lmodern}
\usepackage{microtype}

% ── Checkbox ────────────────────────────────────────────────
\usepackage{amssymb}


% ── Seitenlayout ────────────────────────────────────────────
\usepackage[
  a4paper,
  top=3cm,
  bottom=3cm,
  left=2.5cm,
  right=2.5cm,
  headheight=15pt
]{geometry}

% ── Header & Footer ─────────────────────────────────────────
\usepackage{fancyhdr}
\usepackage{lastpage}

\pagestyle{fancy}
\fancyhf{} % alles zurücksetzen

% Header
\fancyhead[L]{\small\textit{Monster Escape}}
\fancyhead[C]{\small\textit{Informatik 3 – \Semester}}
\fancyhead[R]{\small\textit{\GroupName}}

% Footer
\fancyfoot[L]{\small\AuthorOne \, | \, \AuthorTwo \, | \, \AuthorThree}
\fancyfoot[C]{}
\fancyfoot[R]{\small Seite \thepage\ von \pageref{LastPage}}

% Trennlinie Header/Footer
\renewcommand{\headrulewidth}{0.4pt}
\renewcommand{\footrulewidth}{0.4pt}

% Erste Seite (Titelseite) ohne Header/Footer
\fancypagestyle{plain}{
  \fancyhf{}
  \renewcommand{\headrulewidth}{0pt}
  \renewcommand{\footrulewidth}{0pt}
}

% ── Farben ──────────────────────────────────────────────────
\usepackage[dvipsnames, table]{xcolor}
\definecolor{primaryblue}{RGB}{30, 80, 160}
\definecolor{lightgray}{RGB}{245, 245, 245}
\definecolor{midgray}{RGB}{180, 180, 180}
\definecolor{codebg}{RGB}{248, 248, 248}

% ── Grafiken & Tabellen ─────────────────────────────────────
\usepackage{graphicx}
\usepackage{booktabs}
\usepackage{tabularx}
\usepackage{longtable}
\usepackage{multirow}
\usepackage{array}
\usepackage{colortbl}


% ── Code-Listings ───────────────────────────────────────────
\usepackage{listings}
\lstset{
  backgroundcolor=\color{codebg},
  basicstyle=\ttfamily\small,
  keywordstyle=\color{primaryblue}\bfseries,
  commentstyle=\color{gray}\itshape,
  stringstyle=\color{BrickRed},
  frame=single,
  framerule=0.4pt,
  rulecolor=\color{midgray},
  numbers=left,
  numberstyle=\tiny\color{gray},
  stepnumber=1,
  numbersep=8pt,
  breaklines=true,
  captionpos=b,
  tabsize=2,
  showstringspaces=false
}

% ── Links & PDF-Meta ────────────────────────────────────────
\usepackage{hyperref}
\hypersetup{
  colorlinks=true,
  linkcolor=primaryblue,
  urlcolor=primaryblue,
  citecolor=primaryblue,
  pdftitle={Projektdokumentation - Informatik 3},
  pdfauthor={Gruppenname},
  pdfsubject={2D Spieleprojekt}
}

% ── Sonstige Pakete ─────────────────────────────────────────
\usepackage{enumitem}
\usepackage{titlesec}
\usepackage{tocbibind}
\usepackage{caption}
\usepackage{subcaption}
\usepackage{float}
\usepackage{pifont}   % für \checkmark / \ding
\usepackage{amssymb}

% ── Abschnitts-Formatierung ─────────────────────────────────
\titleformat{\section}
  {\Large\bfseries\color{primaryblue}}
  {\thesection}{1em}{}[\color{primaryblue}\titlerule]

\titleformat{\subsection}
  {\large\bfseries\color{primaryblue!80!black}}
  {\thesubsection}{1em}{}

\titleformat{\subsubsection}
  {\normalsize\bfseries}
  {\thesubsubsection}{1em}{}

% ── Projektdaten (hier anpassen!) ───────────────────────────
\newcommand{\ProjectTitle}{Monster Escape}
\newcommand{\GroupName}{Gruppe B}
\newcommand{\Semester}{SoSe 2026}
\newcommand{\Subject}{Informatik 3}
\newcommand{\Lecturer}{Aysun Cetinalp}
\newcommand{\SubmissionDate}{16.06.2026}
\newcommand{\DocVersion}{1.0}

% Mitglieder
\newcommand{\AuthorOne}{Tobias Hahn}
\newcommand{\AuthorOneMatr}{817616}
\newcommand{\AuthorTwo}{Lutz Morath}
\newcommand{\AuthorTwoMatr}{818145}
\newcommand{\AuthorThree}{Jonas Uzelac}
\newcommand{\AuthorThreeMatr}{804756}

% ============================================================
\begin{document}

% ── Titelseite ──────────────────────────────────────────────
\begin{titlepage}
  \centering
  \vspace*{1.5cm}

  % Obere Linie
  {\color{primaryblue}\rule{\linewidth}{2pt}}
  \vspace{0.5cm}

  {\Huge\bfseries\color{primaryblue} Monster Escape}\\[0.4cm]
  {\Large Projektdokumentation}

  \vspace{0.4cm}
  {\color{primaryblue}\rule{\linewidth}{2pt}}
  \vspace{1.5cm}

  {\large\Subject{} $\cdot$ \Semester}\\[0.3cm]
  {\large Dozentin: Aysun Cetinalp}

  \vfill

% Gruppenmitglieder
  \begin{tabular}{l@{\hspace{30pt}}l}
    \toprule
    \textbf{Name} & \textbf{Matr. Nr.} \\
    \midrule
    \AuthorOne  & \AuthorOneMatr  \\
    \AuthorTwo  & \AuthorTwoMatr  \\
    \AuthorThree & \AuthorThreeMatr \\
    \bottomrule
  \end{tabular}

  \vspace{1cm}

  \begin{tabular}{ll}
    \textbf{Gruppenname:}    & \GroupName      \\
    \textbf{Abgabedatum:}    & \SubmissionDate \\
    % \textbf{Version:}        & \DocVersion     \\
  \end{tabular}

  \vspace{1.5cm}
  {\color{primaryblue}\rule{\linewidth}{1pt}}
\end{titlepage}

% ── Inhaltsverzeichnis ──────────────────────────────────────
\tableofcontents
\newpage

% ── Abbildungsverzeichnis (optional auskommentieren) ────────
% \listoffigures
% \newpage

% ============================================================
\section{Projektbeschreibung}


\begin{description}
  \item[Genre:] Koop 2D Platformer
  \item[Technologie/Engine:] Python + PyGame
  \item[Kurzbeschreibung:] Monster Escape ist ein 2D-Platformer,
  bei dem zwei Spieler gemeinsam an einer Tastatur oder im Multiplayer gemeinsam spielen.
  Ziel ist es, kooperativ einen Ausweg aus dem Level zu finden, indem Rätsel gelöst und
  Hindernisse überwunden werden.
\end{description}

% ============================================================
\section{Spielreferenzen und Inspiration}


\subsection{Fireboy and Watergirl}
\begin{itemize}
  \item \textbf{Genre / Plattform:} Koop 2D Platformer
  \item \textbf{Was wurde übernommen / inspiriert:} \newline Grundspielprinzip, Level Aufbau, Spieldynamik
\end{itemize}

\subsection{Super Mario Bros}
\begin{itemize}
  \item \textbf{Genre / Plattform:} 2D Platformer
  \item \textbf{Was wurde übernommen / inspiriert:} \newline Movement, Grafik
\end{itemize}


\newpage
% ============================================================
\section{Herausforderungen und Lösungsdesign}

\subsection{Grafische Oberfläche}
Um überhaupt ein Window zeichnen zu können und Objekte grafisch darzustellen
haben wir uns dazu entschieden die Bibliothek PyGame zu verwenden, da diese
perfekt zu unserem Anwendungsgebiet passt.

\subsection{Objektplatzierung}
Um Objekte nicht explizit mittels Koordinaten im Code platzieren zu müssen,
nutzen wir den Editor Tiled und die Bibliothek Pytmx. So kann man im Editor
die Objekte grafisch platzieren und mit der entstandenen tmx Datei sowohl Objekte als auch
Grafiken platzieren bzw. darstellen.

\subsection{Inputhandling}
Um eine gute seperation of concerns einzuhalten, wird das komplette Inputhandling in eine
dedizierte InputHandler Klasse ausgelagert. Dort wird der Spieler Input sowie der Input für
für die Menüs verarbeitet.

\subsection{Kollisionsbehandlung}
Da PyGame auch für diese Probleme Lösungen anbietet, verwenden wir die eingebaute Kollisionsbehandlung.
Mit der Klasse pygame.Rect lassen sich Kollisionen leicht überprüfen und weiterverarbeiten

\subsection{Multiplayer}
Hier haben wir uns dazu entschieden, ein Server Client System umzusetzten.
Dafür benutzen wir die Bibliothek socket, um einen TCP Server aufzusetzen.
Ein Spieler ist der Host und startet den Server, der andere muss über die lokale IP Adresse
im gleichen Netzwerk beitreten.

\subsection{Animation}
Für die Animation des Charakters gibt es je nach Bewegungsstatus verschiedene Bilddateien,
durch die Frameweise iteriert wird. So entsteht eine flüssige Animation des Charakters.

\subsection{Verschiedene Menüs}
Um zwischen allen Menüs und Level zu wechseln, nutzen wir einen Game State Manager, welcher
alle Zustände kennt und zwischen ihnen wechselt. Je nach Zustand wird eine andere grafische Oberfläche
dargestellt.

\newpage
% ============================================================
\section{Zeitplan und Aufgabenverteilung}

\begin{table}[H]
  \centering
  \rowcolors{2}{lightgray}{white}
  \begin{tabularx}{\textwidth}{X l l l}
    \toprule
    \rowcolor{primaryblue!20}
    \textbf{Zeitplan} & \textbf{Datum} & \textbf{Verantwortlich} & \textbf{Status} \\
    \midrule
    Projektkonzept \& Anforderungen       & 20.04.2026 & Gemeinsam & $\checkmark$ Erledigt \\
    Kontakt zu Prof.\,Dr.\ Kücherer (EA)  & 20.04.2026 & Gemeinsam & $\checkmark$ Erledigt \\
    UML-Klassendiagramm (Entwurf)         & 20.04.2026 & Gemeinsam & $\checkmark$ Erledigt \\

    Fenster Aufbau                        & 27.04.2026 & Lutz Morath & $\checkmark$ Erledigt \\

    Spieler Physik, Input, Bewegung       & 27.04.2026 & Tobias Hahn & $\checkmark$ Erledigt \\

    Prototyp: Level                       & 04.05.2026 & Jonas Uzelac & $\checkmark$ Erledigt \\

    Kollisionsbehandlung                 & 04.05.2026 & Tobias Hahn & $\checkmark$ Erledigt \\

    Zweiter Spieler lokal                 & 11.05.2026 & Jonas Uzelac & $\checkmark$ Erledigt \\

    Level Mechanik (Start, Ziel, Menu)    & 11.05.2026 & Lutz Morath & $\checkmark$ Erledigt \\

    Tests                                 & 11.05.2026 & Gemeinsam & $\checkmark$ Erledigt \\

    \textbf{Prototyp / spielbare Demo}    & 11.05.2026 & Gemeinsam & $\checkmark$ Erledigt \\

    \textbf{Modellierung Sparx EA fertig} & 01.06.2026 & Gemeinsam & $\checkmark$ Erledigt \\

    Server Client System                  & 20.05.2026 & Tobias Hahn & $\checkmark$ Erledigt \\

    Erweiterte Spielmechaniken            & 02.06.2026 & Jonas Uzelac & $\checkmark$ Erledigt \\

    Tests                                 & 02.06.2026 & Gemeinsam & $\checkmark$ Erledigt \\

    Charakterdesign (Grafik, Animation)   & 02.06.2026 & Lutz Morath & $\checkmark$ Erledigt \\

    Leveldesign (Grafik, Sound)           & 02.06.2026 & Lutz Morath & $\checkmark$ Erledigt \\

    Levelaufbau                           & 02.06.2026 & Jonas Uzelac & $\checkmark$ Erledigt \\


    \textbf{Implementierung abgeschlossen} & 02.06.2026 & Gemeinsam & $\checkmark$ Erledigt \\
    \textbf{Dokumentation abgeschlossen}   & 09.06.2026 & Gemeinsam & $\checkmark$ Erledigt \\
    \textbf{Abgabe}                        & 16.06.2026 & Gemeinsam   & $\checkmark$ Erledigt \\
    \bottomrule
  \end{tabularx}
  \caption{Zeitplan}
\end{table}

\newpage

% ============================================================
\section{UML-Klassendiagramm}

Erstellt mit Sparx Enterprise Architect

\begin{figure}[ht]
    \centering
    % 'page=1' wählt die erste Seite der PDF aus
    % 'width=0.8\textwidth' skaliert das Bild auf 80% der Textbreite
    \includegraphics[page=1, width=1\textwidth]{UML_class.pdf}
    \caption{UML Klassendiagramm dieses Projekts}
    \label{fig:pdf_als_bild}
\end{figure}


% ============================================================
\section{API-Dokumentation}

Kurze Dokumentation der wichtigsten Klassen und ihrer öffentlichen
Schnittstellen.

% Weitere Klassen analog ergänzen

\newpage

% ============================================================
\section{Quellcode-Metriken}

Gemessen zum Zeitpunkt der Abgabe.

\begin{table}[H]
  \centering
  \rowcolors{2}{lightgray}{white}
  \begin{tabular}{ll}
    \toprule
    \rowcolor{primaryblue!20}
    \textbf{Metrik} & \textbf{Wert} \\
    \midrule
    Anzahl Dateien (\#Files)   & 622          \\
    Lines of Code (LOC)        & 1505          \\
    Anzahl Klassen             & 20          \\
    Anzahl Methoden            & 198          \\
    \bottomrule
  \end{tabular}
  \caption{Quellcode-Metriken (ermittelt mit lizard)}
\end{table}

% ============================================================
\section{Retrospektive}

\subsection{Was hat gut funktioniert?}
Die Einbindung des Tiled Map Editors hat gut funktioniert.
Die Level sind jetzt einfach zu gestalten und das Spiel ist somit leicht erweiterbar.

Die Umsetzung der Charakter Animation erwies sich als aufwändiger als gedacht. Das Endprodukt
ist jedoch zufriedenstellend und die Bilddateien für die Animationen sind leicht austauschbar,
wodurch das Aussehen des Spiels einfach verändert werden kann.

Die Spielzustände haben den Arbeitsfluss verbessert. Dadurch war es leicht Menüs zu erweitern oder
neue zu erstellen und neue Level schnell einzubinden. Außerdem wurde dadurch der Umfang des Gameloops
erheblich reduziert.


\subsection{Was würden wir beim nächsten Mal anders machen?}
Obwohl wir mit dem Multiplayer zufrieden sind, gibt es noch Funktionen, die beim Multiplayer nicht
funktionieren. Das Einbinden von beweglichen Objekten in den Multiplayer funktionierte bis zuletzt nicht.
Nächstes mal würden wir bei der Planung der Architektur zu Beginn schon das Server CLient System berücksichtigen.

Außerdem haben wir den zeitlichen Aufwand der Gestaltung unterschätzt. Daher haben wir unseren Spiel Titel
und das Aussehen des Charakters und des Levels geändert, da wir keine passenden Assets für unser
geplantes Thema gefunden haben.


% ============================================================
\end{document}